\chapter{Projectarchitectuur}
\label{app:projectarchitectuur}

\section{Overzicht}
De proof of concept wordt opgevat als een zelfstandige toepassing die
de belangrijkste onderdelen van het aanpassen en verwerken van e-mailtemplates
demonstreert. De frontend, mock-API, templateopslag en renderingpipeline
worden logisch van elkaar gescheiden.

Deze scheiding maakt het mogelijk om de werking van de editor en de
e-mailverwerking te onderzoeken zonder afhankelijk te zijn van de
productieomgeving van Vesalius.ai.

\section{Architectuuroverzicht}

Figuur~\ref{fig:architectuur-overzicht} geeft de beoogde gegevensstroom
schematisch weer.
TODO!!!: nog verdere aanpassingen maken hiervan:
\begin{figure}[htbp]
    \centering
    \begin{verbatim}
+--------------------------------------+
|           Angular-frontend           |
|                                      |
|  Templatekeuze | WYSIWYG-editor      |
|  Preview       | Opslaginterface     |
+------------------+-------------------+
                   |
                   | HTTP/API
                   v
+--------------------------------------+
|       Node.js / Express mock-API     |
|                                      |
|  Validatie | Templatebeheer          |
|  Ophalen   | Opslaan                 |
+------------------+-------------------+
                   |
                   v
+--------------------------------------+
|           Templateopslag             |
|                                      |
|  Templategegevens | Fictieve data    |
|  Tijdelijke opslag of fixtures       |
+------------------+-------------------+
                   |
                   v
+--------------------------------------+
|          Renderingpipeline           |
|                                      |
|  HTML-sanitatie                      |
|  Placeholderverwerking               |
|  Vaste e-mailschil                   |
|  CSS-inlining                        |
|  Eventuele tekstversie               |
+------------------+-------------------+
                   |
                   v
+--------------------------------------+
|          E-mailtestomgeving           |
|                                      |
|              Mailpit                 |
+--------------------------------------+
    \end{verbatim}
    \caption{Beoogde architectuur van de zelfstandige e-maileditor-PoC}
    \label{fig:architectuur-overzicht}
\end{figure}

\noindent
Het schema toont de logische onderdelen en hun onderlinge samenhang. Het is
geen weergave van reeds geverifieerde implementatiedetails.

\section{Frontend}
De frontend vormt de gebruikersinterface(UI) van de toepassing en wordt beoogd
te worden ontwikkeld met Angular en TypeScript.

De frontend is verantwoordelijk voor:
\begin{itemize}
    \item het selecteren en laden van een e-mailtemplate;
    \item het visueel aanpassen van de e-mailinhoud;
    \item het aanroepen van de API om wijzigingen op te slaan;
    \item het opvragen en weergeven van een preview;
    \item het starten van een testmailactie, indien geïmplementeerd.
\end{itemize}

Voor de WYSIWYG-functionaliteit wordt Quill 2 als kandidaat onderzocht.
De definitieve keuze hangt af van een technische proef en de ondersteuning
van de vereiste bewerkingsmogelijkheden.

\section{Mock-API}
De mock-API vormt de verbinding tussen de frontend en de logica voor
templatebeheer. Hiervoor wordt Node.js met TypeScript en Express overwogen.

De API biedt, afhankelijk van de uiteindelijke implementatie, functies voor:
\begin{itemize}
    \item het ophalen van beschikbare templates;
    \item het ophalen van een individuele template;
    \item het valideren en opslaan van wijzigingen;
    \item het aanvragen van een gerenderde preview;
    \item het initiëren van een testmail.
\end{itemize}

De API gebruikt geen productiegegevens of productiegeheimen. Ze dient om
de belangrijkste interacties van de PoC te demonstreren.

\section{Templateopslag en gegevensmodel}
De opslagmethode wordt gekozen op basis van de behoeften van de PoC.
Een tijdelijke opslagmethode, zoals JSON-fixtures of een verwisselbare
opslagimplementatie, kan volstaan zolang de vereiste functionaliteiten
aantoonbaar blijven.

De architectuur moet vermijden dat de volledige toepassingslogica afhankelijk
wordt van één specifieke opslagtechnologie. Daardoor kan een mogelijke
latere integratie een andere opslaglaag gebruiken zonder de volledige
rendering- en editorlogica te moeten herschrijven.

Een conceptueel template bevat de volgende gegevens:

\begin{table}[htbp]
    \centering
    \caption{Conceptuele templatevelden}
    \label{tab:templatevelden}
    \begin{tabular}{|p{3.5cm}|p{9cm}|}
        \hline
        \textbf{Veld} & \textbf{Beschrijving} \\
        \hline
        Identificatie & Unieke identificatie van de template. \\
        \hline
        E-mailtype & Type e-mail, bijvoorbeeld bevestiging of herinnering. \\
        \hline
        Taal & Taal van de e-mailinhoud, indien ondersteund. \\
        \hline
        Inhoud & Bewerkbare inhoud van de template. \\
        \hline
        Placeholders & Verwijzingen naar toegestane dynamische gegevens. \\
        \hline
        Configuratie & Eventuele configuratie die nodig is om de template
        correct te renderen. \\
        \hline
        Status & Eventuele status of versie-informatie, indien voorzien. \\
        \hline
    \end{tabular}
\end{table}

Dit overzicht is een conceptueel model. De uiteindelijke velden, types en
validatieregels moeten overeenstemmen met het geïmplementeerde gegevensmodel.

\section{Renderingpipeline}
De renderingpipeline zet de bewerkbare template om in een e-mail die geschikt
is voor de testomgeving. De precieze volgorde en implementatie worden
vastgelegd op basis van de gerealiseerde oplossing.

De belangrijkste verwerkingsstappen zijn:

\begin{enumerate}
    \item \textbf{Validatie:} controleer of de ontvangen template en de
    aangeleverde gegevens aan de afgesproken voorwaarden voldoen.
    \item \textbf{HTML-sanitatie:} verwijder of weiger niet-toegestane
    HTML-elementen en attributen volgens een expliciete allowlist.
    \item \textbf{Placeholderverwerking:} vervang alleen toegestane
    placeholders door de bijbehorende fictieve waarden.
    \item \textbf{Samenstellen van de e-mail:} combineer de inhoud met de
    vaste e-mailschil, indien die deel uitmaakt van de implementatie.
    \item \textbf{CSS-verwerking:} plaats ondersteunde CSS-regels inline
    waar dat nodig is voor compatibiliteit met e-mailclients.
    \item \textbf{Tekstversie:} genereer indien voorzien een alternatieve
    tekstversie van de e-mail.
    \item \textbf{Uitvoer:} geef de gerenderde inhoud terug voor de preview
    of bied ze aan de lokale testmailservice aan.
\end{enumerate}

De verwerking moet rekening houden met de context waarin placeholders
worden gebruikt. HTML-escaping, toegestane waarden en veilige verwerking
van invoer blijven noodzakelijk, ook wanneer de template zelf wordt
gesanitiseerd.

\section{E-mailschil en piloottemplates}
De PoC richt zich op drie e-mailtypes:
\begin{itemize}
    \item afspraakbevestiging;
    \item afspraakherinnering;
    \item afspraakannulering.
\end{itemize}

De vaste e-mailschil kan onder meer het logo, de merkkleuren en de footer
bevatten. De nadruk ligt op het visueel aanpassen van de e-mailinhoud en
niet op een volledige editor voor de gehele e-maillay-out.

\section{Testinfrastructuur}
Voor de lokale testinfrastructuur worden Docker en Mailpit overwogen.
Mailpit kan dienen om testmails lokaal op te vangen en te inspecteren zonder
dat daarvoor echte ontvangers nodig zijn.

Geautomatiseerde tests kunnen worden uitgevoerd met Vitest en Playwright,
afhankelijk van de uiteindelijke teststrategie.

De tests richten zich op:
\begin{itemize}
    \item templatevalidatie en opslag;
    \item placeholderverwerking;
    \item HTML-sanitatie;
    \item rendering en preview;
    \item de belangrijkste gebruikersflows;
    \item het versturen en controleren van testmails.
\end{itemize}

\section{Privacy en beveiliging}
De PoC gebruikt uitsluitend fictieve gegevens. Er worden geen echte
patiëntgegevens verwerkt en er worden geen productiecredentials in de
broncode of testbestanden opgenomen.

De belangrijkste beveiligingsmaatregelen zijn:
\begin{itemize}
    \item een expliciete allowlist voor toegestane HTML-inhoud;
    \item validatie en veilige verwerking van placeholders;
    \item gescheiden configuratie van test- en productieomgevingen;
    \item het vermijden van echte persoonsgegevens en geheimen;
    \item het lokaal opvangen van testmails.
\end{itemize}

Deze maatregelen verminderen bepaalde risico's, maar maken de PoC niet
automatisch geschikt voor productiegebruik.

\section{Mogelijke latere integratie}
Een mogelijke integratie in Vesalius.ai valt buiten de scope van dit project.
De PoC kan wel informatie opleveren over de benodigde API-contracten, de
opslaginterface, de renderingpipeline, de veiligheidsvereisten en de
compatibiliteit met e-mailclients.

Een integratievoorstel moet daarom rekening houden met de bestaande
architectuur van Vesalius.ai, de interne ontwikkelrichtlijnen en de
relevante privacy- en beveiligingsvereisten.